%================================================================
\chapter{Problem Formulation and Dataset}
\label{ch:problem}
%================================================================

\section{Network Graph Model}
\label{sec:graph_model}

A water distribution network is modelled as a directed graph
$\mathcal{G} = (\mathcal{V}, \mathcal{E})$. The node set
$\mathcal{V}$ contains all junctions, tanks, and reservoirs. The
edge set $\mathcal{E}$ contains all pipes, pumps, and valves, each
directed by the convention that positive flow runs from a lower
node index to a higher one (the direction is reversed in the feature
vector if the actual flow is negative).

\subsection{Node Features}

Each node $i \in \mathcal{V}$ carries a feature vector
$\mathbf{x}_i \in \mathbb{R}^7$:
\begin{enumerate}
  \item Observed pressure $\tilde{p}_i$ (metres; zero if missing);
  \item Elevation (metres above sea level);
  \item Base demand ($\mathrm{m}^3/\mathrm{s}$);
  \item Demand pattern scaling factor at the current time step;
  \item Observation mask $m_i^p \in \{0,1\}$ (1 if pressure observed);
  \item Node type, one-hot over \{junction, tank, reservoir\}: two bits.
\end{enumerate}
Continuous features are normalised to zero mean and unit standard
deviation using statistics computed on the training split.

\subsection{Edge Features}

Each pipe $(i,j) \in \mathcal{E}$ carries a feature vector
$\mathbf{e}_{ij} \in \mathbb{R}^8$:
\begin{enumerate}
  \item Observed flow $\tilde{q}_{ij}$ (zero if missing);
  \item Pipe length (metres);
  \item Pipe diameter (metres);
  \item Roughness coefficient (Hazen-Williams $C$);
  \item Observation mask $m_{ij}^q$;
  \item Pipe type, one-hot over \{pipe, pump, valve\}: three bits.
\end{enumerate}

\subsection{Graph Snapshot}

A \emph{snapshot} is one time-step observation of the full network.
At time step $t$, the snapshot is a heterogeneous attributed graph
\begin{equation}
  s_t = \bigl(\mathbf{X}_t, \mathbf{E}_t, \mathcal{G}\bigr),
\end{equation}
where $\mathbf{X}_t \in \mathbb{R}^{N \times 7}$ is the node feature
matrix and $\mathbf{E}_t \in \mathbb{R}^{M \times 8}$ is the edge
feature matrix. The graph topology $\mathcal{G}$ is fixed across
all time steps (the network does not change; only the readings
change).

A \emph{temporal window} of length $T$ is the sequence
$w_t = (s_{t-T+1}, \ldots, s_t)$. All temporal models in this
thesis use $T = 6$.

\section{Task Definition}
\label{sec:task}

\subsection{Reconstruction}

Given a window $w_t$ with missing sensor values, predict:
\begin{equation}
  \hat{p}_i \approx p_i^*, \quad \forall i \in \mathcal{V},
  \qquad
  \hat{q}_{ij} \approx q_{ij}^*, \quad \forall (i,j) \in \mathcal{E},
\end{equation}
where $p_i^*$ and $q_{ij}^*$ are the ground-truth values from the
hydraulic simulation. Performance is measured by MAE
and RMSE, reported separately for observed sensors
(which the model has access to) and unobserved sensors (which it
must impute).

\subsection{Anomaly Detection}

Given the window $w_t$ (which may contain corrupted readings),
predict a binary anomaly label at each sensor:
\begin{equation}
  \hat{a}_i \in [0,1], \quad \forall i \in \mathcal{V},
  \qquad
  \hat{a}_{ij} \in [0,1], \quad \forall (i,j) \in \mathcal{E}.
\end{equation}
A threshold of 0.5 is used to convert the logit to a binary
prediction. Performance is measured by precision, recall, F1 score,
and AUROC (threshold-free). The primary metric throughout
this thesis is the pressure anomaly F1 score, which is
reported per-attack and as an aggregate over all attack types.

\section{Attack Taxonomy}
\label{sec:attacks}

\subsection{Attack Model}

An attacker can modify the observed pressure reading $\tilde{p}_i$
at any sensor $i$. The attacker has a budget: it can touch at most
$k$ sensors per snapshot, and the magnitude of any single
perturbation is bounded by $\varepsilon$ in the $L_\infty$ norm.
Formally:
\begin{equation}
  \|\boldsymbol{\delta}_p\|_\infty \leq \varepsilon, \qquad
  \|\boldsymbol{\delta}_p\|_0 \leq k,
\end{equation}
where $\boldsymbol{\delta}_p \in \mathbb{R}^N$ is the perturbation
vector. The same constraints apply to flow sensors.

\subsection{Five Hand-Crafted Attack Types}

The corruption pipeline generates five attack types alongside clean
snapshots. During training the attacker type is known (used to
supervise the router); at test time the model receives only the
corrupted reading.

\textbf{Random falsification (type 1).} For each attacked sensor $i$
the corrupted reading is:
\begin{equation}
  \tilde{p}_i = \alpha_i \, p_i^* + \beta_i,
\end{equation}
where $\alpha_i \sim \mathcal{U}(0.8, 1.2)$ and $\beta_i \sim
\mathcal{U}(-3, 3)$ are drawn independently. This represents an
attacker who does not know or does not care about the hydraulic
state.

\textbf{Replay attack (type 2).} A lag $\tau \sim \mathcal{U}(1, T)$
is sampled once per snapshot. For each attacked sensor:
\begin{equation}
  \tilde{p}_i(t) = p_i^*(t - \tau).
\end{equation}
The replayed value is a genuinely observed past measurement, so it
lies on the data manifold. This is the hardest attack for spatial
models.

\textbf{Stealthy bias injection (type 3).} A slow, linear drift is
applied over the window:
\begin{equation}
  \tilde{p}_i(t) = p_i^*(t) + \mu \cdot \frac{t - t_0}{\Delta t},
\end{equation}
where $\mu \sim \mathcal{U}(0.5, 2.0)$ is the drift rate and $\Delta t$
is the number of steps since the attack started. The attack is
designed to be hard to detect at any single time step because the
anomaly magnitude is initially small.

\textbf{Noise injection (type 4).} High-variance Gaussian noise is
added:
\begin{equation}
  \tilde{p}_i = p_i^* + \epsilon, \quad
  \epsilon \sim \mathcal{N}(0, \sigma_\text{atk}^2),
\end{equation}
where $\sigma_\text{atk}$ is set to five times the background sensor
noise level. This type is easy to detect but serves as a baseline
comparison.

\textbf{Targeted attack (type 5).} Sensors are selected with
probability proportional to their degree-weighted demand. The
perturbation is the same as random falsification but applied to
high-impact sensors:
\begin{equation}
  P(\text{attack } i) \propto d_i \cdot \deg(i).
\end{equation}
In the spatial \gnn{}, high-degree nodes tend to have
higher-quality embeddings because they aggregate more neighbourhood
information, making targeted attacks somewhat harder to detect than
purely random ones.

Table~\ref{tab:attacks} summarises the five attack types.

\begin{table}[htbp]
  \centering
  \caption{Summary of the five hand-crafted attack types. The
  detection difficulty rating is qualitative and refers to the
  pretrained Temporal \moe{} defender on Modena.}
  \label{tab:attacks}
  \begin{tabular}{llll}
    \toprule
    Type & Name & Key characteristic & Detection (F1) \\
    \midrule
    1 & Random     & Scale + bias         & 0.996 \\
    2 & Replay     & Past value broadcast & 0.196 \\
    3 & Stealthy   & Slow drift           & 0.927 \\
    4 & Noise      & High variance        & 0.902 \\
    5 & Targeted   & High-impact sensors  & 0.975 \\
    \bottomrule
  \end{tabular}
\end{table}

\subsection{Two Novel Held-Out Attack Types}
\label{sec:heldout_attacks}

To test generalisation beyond the training distribution, two attack
types are synthesised that none of the hand-crafted training attacks
resemble.

\textbf{Sinusoidal injection.} A sinusoidal signal is added to the
true reading:
\begin{equation}
  \tilde{p}_i(t) = p_i^*(t) + A \sin\!\left(
    \frac{2\pi t}{T_\text{period}} + \phi
  \right),
\end{equation}
where $A \sim \mathcal{U}(1, 4)$, $T_\text{period} \in \{3, 4, 6\}$
steps, and $\phi \sim \mathcal{U}(0, 2\pi)$. None of the five
training types produces a periodic signal.

\textbf{Cross-sensor swap.} The readings of two attacked sensors are
exchanged:
\begin{equation}
  \tilde{p}_i = p_j^*, \qquad \tilde{p}_j = p_i^*.
\end{equation}
A swap is locally plausible at each sensor individually (the swapped
reading was a real pressure at another sensor) but globally
inconsistent because the two sensors have different neighbourhoods.

\section{Dataset Generation}
\label{sec:dataset}

\subsection{Simulation Setup}

Each dataset is generated by running the WNTR hydraulic
simulator on the corresponding EPANET input file for 24
hours at a one-hour time step, repeated over 50 demand scenarios.
In each scenario the base demands at all junctions are perturbed by
$\pm 20\%$ using a uniform random factor, and a sinusoidal diurnal
pattern is applied on top. This produces 50 independent sequences
of 24 snapshots each, for a total of 1,200 snapshots per network
before corruption.

\subsection{Missing Data}

After simulation, 50\% of pressure sensors and 50\% of flow sensors
are randomly removed per snapshot. The same rate is used for all
networks and all experiments to ensure comparability.

\subsection{Attack Injection}

For each snapshot, one of the five attack types (or the clean class)
is selected uniformly at random. The attack fraction parameter
$f_\text{atk} = 0.15$ determines the expected fraction of sensors
attacked per snapshot: sensors are independently attacked with this
probability. In practice, the mean number of attacked sensors is
$0.15 \times N$ for pressure and $0.15 \times M$ for flow.

\subsection{Train / Validation / Test Split}

The 50 scenarios are split by scenario index: 35 training, 8
validation, 7 test. This means all time steps within a scenario
appear in the same split, preventing temporal leakage. The temporal
model uses a sliding window of $T = 6$ steps, so the first five
steps of each scenario are excluded to avoid partial-window inputs.

Table~\ref{tab:dataset_stats} shows dataset sizes for each network.

\begin{table}[htbp]
  \centering
  \caption{Dataset sizes after windowing ($T = 6$, attack fraction 0.15).}
  \label{tab:dataset_stats}
  \begin{tabular}{lrrrr}
    \toprule
    Network & Train windows & Val windows & Test windows & Model params \\
    \midrule
    Net1    & 532 & 122 & 105 & $\sim$32k \\
    Net3    & 532 & 122 & 105 & $\sim$219k \\
    Modena  & 532 & 122 & 105 & $\sim$219k \\
    \bottomrule
  \end{tabular}
\end{table}

\subsection{Evaluation Protocol}

All results are reported on the held-out test split. Models are
selected by the best validation anomaly F1 during training, not by
reconstruction loss, because reconstruction-optimal checkpoints
tend to be replay-blind: they over-fit to anomaly-free windows.
For the self-play experiments, the selection criterion is a composite
of anomaly F1 and replay F1 to prevent checkpoints that improve
overall F1 by sacrificing replay recall entirely.

\section{Classical Baselines}
\label{sec:classical_baselines}

Before evaluating learned models, three classical reconstruction
methods establish the difficulty of the task.

\textbf{Mean imputation} replaces each missing value with the
per-sensor training mean. This completely ignores graph structure
and provides a lower bound on reconstruction quality.

\textbf{Pseudo-inverse.} The pressure vector is recovered by solving
the under-determined system $\mathbf{A}_o \mathbf{p} = \mathbf{y}$,
where $\mathbf{A}_o$ is the sub-matrix of the incidence matrix
restricted to observed sensors, using the Moore-Penrose pseudo-inverse.

\textbf{Weighted Least Squares (WLS).} The same linear system is
solved with weights inversely proportional to sensor noise variance.
This is the standard approach in hydraulic state estimation.

For anomaly detection, all three methods use a threshold on the
reconstruction residual: a sensor is flagged anomalous if its
residual exceeds three standard deviations of the clean-data residual
distribution. The AUROC is computed by varying this
threshold.

The results in Section~\ref{sec:classical_results} show that classical
methods achieve anomaly F1 scores below 0.15 on all three networks
and reconstruction MAE values 10-50 times worse than the
\gnn{}-based models.
